% Originally: content/week-09.tex, \section{Solutions}

\section{Solutions}
\label{sec:gram_determinant_solutions}

\begin{exercisesolution}[Solution to \cref{ex:area_formula_two_maps}]
\label{sol:area_formula_two_maps}
% Generator: Claude Opus 5 (High) --- the five statements are the examiner's; the working, the
% counterexample and the AI-Note below are written for these notes. The verdicts (a) true,
% (b)-(d) false, (e) true are the official key's; ours agree on (b)-(e) and diverge on (a), for
% the reason set out in the AI-Note.
Take the second map first, because there \eqref{eq:mock_area_formula} applies without any
difficulty. Write $\Psi = L + c$ with
\[L := \begin{pmatrix} 1 & 1 \\ \alpha & -1 \\ 0 & 1 \end{pmatrix}, \qquad
  c := \begin{pmatrix} 0 \\ 0 \\ 1 \end{pmatrix},\]
so $\Jac\Psi(x,y) = L$ at every point. The Gram matrix of its two columns
$c_1 = (1,\alpha,0)\transp$ and $c_2 = (1,-1,1)\transp$ is
\[L\transp L = \begin{pmatrix} \langle c_1,c_1\rangle & \langle c_1,c_2\rangle \\
  \langle c_2,c_1\rangle & \langle c_2,c_2\rangle \end{pmatrix}
  = \begin{pmatrix} 1+\alpha^2 & 1-\alpha \\ 1-\alpha & 3 \end{pmatrix},\]
whose determinant is
\[3(1+\alpha^2) - (1-\alpha)^2 = 3 + 3\alpha^2 - 1 + 2\alpha - \alpha^2 = 2(\alpha^2+\alpha+1).\]
The integrand of \eqref{eq:mock_area_formula} is therefore the constant
$\sqrt{2(\alpha^2+\alpha+1)}$, and integrating a constant over $E$ gives
\[\vol_2(\Psi(E)) = \sqrt{2(\alpha^2+\alpha+1)}\ \vol_2(E).\]
So \textbf{(e)} is true and \textbf{(c)} and \textbf{(d)} are false. Note that
$\alpha^2+\alpha+1$ has discriminant $-3$ and so is strictly positive for every real $\alpha$:
the two columns are independent whatever $\alpha$ is, $\Psi$ is an injective affine map, and the
formula is exact. Statement \textbf{(c)} is what an $\alpha$-linear guess produces, and
\textbf{(d)} is what one gets by taking a square root before, rather than after, forming the
Gram determinant.

For $\Phi$ the Jacobian is square, and there $\sqrt{\det(\Jac\Phi\transp\Jac\Phi)} =
\lvert\det\Jac\Phi\rvert$. Since
\[\Jac\Phi(x,y) = \begin{pmatrix} y & x \\ 2x & -2y \end{pmatrix}, \qquad
  \det\Jac\Phi(x,y) = -2y^2 - 2x^2,\]
the integrand is $2(x^2+y^2)$. Applied literally, \eqref{eq:mock_area_formula} therefore returns
$2\int_E(x^2+y^2)$, which is statement \textbf{(a)}, and never $2\int_E(x^2+2y^2)$, so
\textbf{(b)} is false.
\end{exercisesolution}

\begin{ainote}[The quoted area formula needs injectivity, and $\Phi$ does not have it]
\label{note:area_formula_needs_injectivity}
\Cref{ex:area_formula_two_maps}\textbf{(a)} is marked \emph{true} in the official key, and it is
what \eqref{eq:mock_area_formula} gives. But \eqref{eq:mock_area_formula} is false as printed,
because it carries no injectivity hypothesis, and $\Phi$ is exactly the map that exposes this:
$\Phi(-p) = \Phi(p)$ for every $p$.

Concretely, take $E := \overline{B_1(0)}$ and put $u := xy$, $v := x^2-y^2$. Then
\[4u^2 + v^2 = 4x^2y^2 + (x^2-y^2)^2 = (x^2+y^2)^2 ,\]
so $\Phi$ maps $E$ into the ellipse $\{4u^2+v^2 \leq 1\}$, and onto it: given such a $(u,v)$,
set $s := \sqrt{4u^2+v^2} \leq 1$ and solve $x^2 = (s+v)/2$, $y^2 = (s-v)/2$, both non-negative
because $v^2 \leq s^2$. These satisfy $x^2+y^2 = s \leq 1$ and $(xy)^2 = (s^2-v^2)/4 = u^2$, so a
choice of signs gives $xy = u$. Hence
\[\Phi(E) = \{(u,v) \mid 4u^2+v^2 \leq 1\},\]
an ellipse with semi-axes $\tfrac12$ and $1$, of area $\pi/2$. The right-hand side of
\eqref{eq:mock_area_formula}, meanwhile, is
\[\int_E 2(x^2+y^2)\,dx\,dy = 2\int_0^{2\pi}\!\!\int_0^1 r^2\cdot r\,dr\,d\theta = \pi .\]
The two differ by a factor of $2$, which is precisely the number of preimages. What is true in
general is that the integral counts multiplicity,
\[\int_E \lvert\det\Jac\Phi\rvert\,dx = \int_{\mathbb{R}^n} \#\big(\Phi^{-1}(w)\cap E\big)\,dw
  \ \geq\ \vol_n(\Phi(E)),\]
with equality exactly when $\Phi$ is injective off a null set. \Cref{thm:change_of_variables}
carries the injectivity hypothesis for this reason, and so does
\cref{def:d_volume_parametrized_submanifold}, which is stated for a \emph{parametrization} rather
than for an arbitrary $C^1$ map.

The verdict on part \textbf{(a)} is therefore: true as the examination intends it, false as
mathematics. It is kept here because the gap is more instructive than the computation.
\end{ainote}

\begin{exercisesolution}[Solution to \cref{ex:paraboloid_hypersurface_volume}]
\label{sol:paraboloid_hypersurface_volume}
% This label is referenced from 02-volume-of-embedded-surfaces.tex with \cpageref, NOT \cref.
% exercisesolution wraps `proof`, which has no counter, so a \label inside it is silently
% misattributed to the last-stepped ambient counter: \cref printed "Section 21.e" here until
% 2026-08-09. Page references are unaffected and are the documented workaround in style.md.
\label{ex:solution_9.1}
% Generator: Claude Sonnet 5 (Medium)
Following the official hint rather than Corsin's (\cref{note:9.1_coordinates_mislabelled}), we use
cylindrical coordinates $(x,y,z) = (r\cos\theta,r\sin\theta,z)$, under which $x^2+y^2=r^2$.
The two bounding surfaces become $r^2+z^2=8$ and $z = r^2/2$; intersecting them,
\[z^2+2z-8 = 0 \quad\Longrightarrow\quad z = 2 \text{ or } z=-4,\]
and since $z=r^2/2\geq0$ on the paraboloid, the surfaces meet at $z=2$, $r=2$. For
$r \in (0,2)$ the paraboloid $z=r^2/2$ lies below the upper half of the sphere
$z=\sqrt{8-r^2}$, so
\[B = \left\{(r,\theta,z) : 0<r<2,\ 0<\theta<2\pi,\ \tfrac12r^2 < z < \sqrt{8-r^2}\right\}.\]
Hence, with $dx\,dy\,dz = r\,dr\,d\theta\,dz$,
\[
\begin{aligned}
\vol(B) &= \int_0^{2\pi}\!\!\int_0^2 r\left(\sqrt{8-r^2} - \frac{r^2}{2}\right)dr\,d\theta
  = 2\pi\int_0^2 \left(r\sqrt{8-r^2} - \frac{r^3}{2}\right)dr.
\end{aligned}
\]
For the first term, substitute $u = 8-r^2$, $du=-2r\,dr$:
\[\int_0^2 r\sqrt{8-r^2}\,dr = \left[-\frac{u^{3/2}}{3}\right]_{u=8}^{u=4}
  = \frac{8^{3/2}-4^{3/2}}{3} = \frac{16\sqrt2-8}{3}.\]
For the second, $\int_0^2 \tfrac12 r^3\,dr = \tfrac12\cdot\tfrac{16}{4} = 2$. So
\[\vol(B) = 2\pi\left(\frac{16\sqrt2-8}{3} - 2\right) = 2\pi\cdot\frac{16\sqrt2-14}{3}
  = \boxed{\ \frac{4\pi}{3}\big(8\sqrt2-7\big)\ }.\]
\end{exercisesolution}

\begin{exercisesolution}[Solution to \cref{ex:curve_length_polar_and_vector_integral}]
\label{sol:curve_length_polar_and_vector_integral}
% Generator: Claude Sonnet 5 (Medium)
\begin{enumerate}[label=\textbf{(\alph*)}]
  \item By Fubini (\cref{thm:fubini}),
    \[\iint_D (x^3+3x^2y+y^3)\,dx\,dy = \int_0^2\int_0^1 (x^3+3x^2y+y^3)\,dy\,dx
      = \int_0^2\left(x^3+\frac32x^2+\frac14\right)dx.\]
    The last integral is
    $\left[\tfrac{x^4}{4}+\tfrac{x^3}{2}+\tfrac{x}{4}\right]_0^2 = 4+4+\tfrac12
    = \boxed{\ \tfrac{17}{2}\ }$.
  \item Using Corsin's parametrization $D=\{0\leq x\leq y\leq\pi\}$,
    \[\iint_D x\cos(x+y)\,dx\,dy = \int_0^\pi\int_x^\pi x\cos(x+y)\,dy\,dx
      = \int_0^\pi x\big[\sin(x+y)\big]_{y=x}^{y=\pi}\,dx.\]
    Since $\sin(x+\pi) = -\sin x$, the bracket equals $-\sin x - \sin(2x)$, so
    \[\iint_D x\cos(x+y)\,dx\,dy = -\int_0^\pi x\sin x\,dx - \int_0^\pi x\sin(2x)\,dx.\]
    Integrating by parts, $\int_0^\pi x\sin x\,dx = [-x\cos x+\sin x]_0^\pi = \pi$, and
    $\int_0^\pi x\sin(2x)\,dx = \left[-\tfrac{x}{2}\cos(2x)+\tfrac14\sin(2x)\right]_0^\pi
    = -\tfrac\pi2$. Hence
    \[\iint_D x\cos(x+y)\,dx\,dy = -\pi - \left(-\frac\pi2\right) = \boxed{\ -\frac\pi2\ }.\]
  \item Here $D = \{1<x<2,\ 1<y<3-x\}$, so
    \[\iint_D \frac{1}{(x+y)^3}\,dx\,dy = \int_1^2\int_1^{3-x} (x+y)^{-3}\,dy\,dx
      = \int_1^2 \left[-\frac12(x+y)^{-2}\right]_{y=1}^{y=3-x}dx.\]
    Since $x+(3-x)=3$, the bracket is $-\tfrac12\big(3^{-2}-(x+1)^{-2}\big)
    = \tfrac12(x+1)^{-2} - \tfrac1{18}$, and
    \[\int_1^2\left(\frac{1}{2(x+1)^2}-\frac1{18}\right)dx
      = \left[-\frac{1}{2(x+1)}\right]_1^2 - \frac1{18}
      = \left(-\frac16+\frac14\right) - \frac1{18} = \frac1{12}-\frac1{18}
      = \boxed{\ \frac1{36}\ }.\]
\end{enumerate}
\end{exercisesolution}

% All three parts of ex:curve_length_polar_and_vector_integral and ex:paraboloid_hypersurface_volume were worked out independently before consulting
% Sol9_Analysis2_eng.pdf; all four answers (4pi/3 (8sqrt2 - 7), 17/2, -pi/2, 1/36) match the
% official solution exactly, so there is no divergence to record. (Was an ainote until
% 2026-08-09: style.md names cross-checking against SolN as editor-facing under Test 2.)

\begin{exercisesolution}[Solution to \cref{ex:spherical_quadrilateral_area}]
\label{sol:spherical_quadrilateral_area}
% Generator: Claude Sonnet 5 (Medium)
We parametrize $\mathbb{S}^2$ (up to a null set) using latitude/longitude spherical
coordinates
\[\Phi : (-\pi,\pi)\times\left(-\tfrac\pi2,\tfrac\pi2\right) \to \mathbb{S}^2, \qquad
  (\varphi,\theta) \mapsto \begin{pmatrix}\cos\varphi\cos\theta\\ \sin\varphi\cos\theta\\
  \sin\theta\end{pmatrix}.\]
Note that this is \emph{not} the $(r\equiv1)$ parametrization used elsewhere in this chapter,
whose polar angle runs over $(0,\pi)$. Here $\theta$ is measured from the equator, matching
the exercise's range $\theta\in(-\pi/2,\pi/2)$. Computing
\[\partial_\varphi\Phi\times\partial_\theta\Phi = \begin{pmatrix}-\sin\varphi\cos\theta\\
  \cos\varphi\cos\theta\\ 0\end{pmatrix} \times \begin{pmatrix}-\cos\varphi\sin\theta\\
  -\sin\varphi\sin\theta\\ \cos\theta\end{pmatrix} = \begin{pmatrix}\cos\varphi\cos^2\theta\\
  \sin\varphi\cos^2\theta\\ \sin\theta\cos\theta\end{pmatrix},\]
its length is $\lvert\partial_\varphi\Phi\times\partial_\theta\Phi\rvert = \cos\theta$ (for
$\theta\in(-\pi/2,\pi/2)$, so $\cos\theta>0$). By \cref{prop:flux_r3}\textbf{(b)}, this is the
Gram determinant, so $dA = \cos\theta\,d\theta\,d\varphi$, and
\[\int_S dA = \int_{\varphi_1}^{\varphi_2}\!\!\int_{\theta_1}^{\theta_2} \cos\theta\,d\theta\,
  d\varphi = \boxed{\ (\varphi_2-\varphi_1)\big(\sin\theta_2-\sin\theta_1\big)\ }.\]
\end{exercisesolution}

\begin{exercisesolution}[Solution to \cref{ex:torus_surface_area}]
\label{sol:torus_surface_area}
% Generator: Claude Sonnet 5 (Medium)
The circle $S$ has radius $2$, and $T$ is the solid torus swept by a tube of radius $1$
around it, so $\partial T$ is Corsin's general torus surface with $R=2$, $r=1$ fixed. Using
his parametrization,
\[\phi(\theta,\varphi) = \begin{pmatrix}(2+\cos\varphi)\cos\theta\\
  (2+\cos\varphi)\sin\theta\\ \sin\varphi\end{pmatrix}, \qquad \theta,\varphi\in(0,2\pi).\]
Then $\partial_\theta\phi = \big(-(2+\cos\varphi)\sin\theta,\ (2+\cos\varphi)\cos\theta,\
0\big)$ has length $2+\cos\varphi$, and $\partial_\varphi\phi = (-\sin\varphi\cos\theta,\
-\sin\varphi\sin\theta,\ \cos\varphi)$ has length $1$; the two are orthogonal (their dot
product is $\sin\varphi\cos\varphi\sin\theta\cos\theta - \sin\varphi\cos\varphi\sin\theta
\cos\theta = 0$), so
\[\sqrt{\det D\phi\transp D\phi} = \lvert\partial_\theta\phi\rvert\,\lvert\partial_\varphi\phi
  \rvert = 2+\cos\varphi.\]
Hence
\[\vol_2(\partial T) = \int_0^{2\pi}\!\!\int_0^{2\pi} (2+\cos\varphi)\,d\theta\,d\varphi
  = 2\pi\left(2\cdot2\pi + \underbrace{\int_0^{2\pi}\cos\varphi\,d\varphi}_{=0}\right)
  = \boxed{\ 8\pi^2\ }.\]
\end{exercisesolution}

\begin{exercisesolution}[Solution to \cref{ex:solids_of_revolution_area}]
\label{sol:solids_of_revolution_area}
% Generator: Claude Sonnet 5 (Medium)
\begin{enumerate}[label=\textbf{(\alph*)}]
  \item $\partial_{\mathrm{side}}\Omega$ is not closed (it omits, e.g., the limit
    point $(a,f(a),0)$ as $x\to a^+$) and not open (a $2$-dimensional surface has empty interior
    in $\mathbb{R}^3$), but it \emph{is} path-connected: from any point, first move along the
    graph of $f$ in the $x$-direction, then move around the circle in the $(y,z)$-variables.
    $\Omega$ is open (as the strict sublevel set of a continuous function), not closed, and its
    connectedness depends on $f$: if $f$ vanishes somewhere in $(a,b)$, the cross-section pinches
    to a point there and $\Omega$ is disconnected into the pieces on either side; if $f>0$
    throughout $(a,b)$, $\Omega$ is connected.

  \item In cylindrical coordinates $(x,\rho,\theta)$ with $y=\rho\cos\theta$,
    $z=\rho\sin\theta$, $dx\,dy\,dz = \rho\,d\rho\,d\theta\,dx$, the domain becomes the box-like
    region $\{a<x<b,\ 0<\rho<f(x),\ \theta\in[0,2\pi)\}$, so
    \[\vol_3(\Omega) = \int_a^b\!\!\int_0^{2\pi}\!\!\int_0^{f(x)} \rho\,d\rho\,d\theta\,dx
      = \int_a^b\!\!\int_0^{2\pi} \frac12 f(x)^2\,d\theta\,dx
      = \boxed{\ \pi\int_a^b f(x)^2\,dx\ }.\]

  \item Using the suggested parametrization $\Phi(x,\theta) = (x,f(x)\cos\theta,
    f(x)\sin\theta)$, write the orthonormal frame $\xi_1:=(1,0,0)$, $\xi_2:=(0,\cos\theta,
    \sin\theta)$, $\xi_3:=(0,-\sin\theta,\cos\theta)$, so that
    \[\partial_x\Phi = \xi_1 + f'(x)\xi_2, \qquad \partial_\theta\Phi = f(x)\xi_3.\]
    Since $\xi_1,\xi_2,\xi_3$ are orthonormal and right-handed, $\xi_1\times\xi_3=-\xi_2$ and
    $\xi_2\times\xi_3=\xi_1$, so
    \begin{align*}
      \partial_x\Phi\times\partial_\theta\Phi
        &= f(x)\big(\xi_1\times\xi_3 + f'(x)\,\xi_2\times\xi_3\big)
         = f(x)\big({-\xi_2}+f'(x)\xi_1\big), \\
      \lvert\partial_x\Phi\times\partial_\theta\Phi\rvert &= f(x)\sqrt{1+f'(x)^2}.
    \end{align*}
    By \cref{prop:flux_r3}\textbf{(b)}, this is the Gram determinant, so
    \[\vol_2(\partial_{\mathrm{side}}\Omega) = \int_a^b\!\!\int_0^{2\pi} f(x)\sqrt{1+f'(x)^2}\,
      d\theta\,dx = \boxed{\ 2\pi\int_a^b f(x)\sqrt{1+f'(x)^2}\,dx\ }.\]

  \item With $f(x)=1/x$ on $[1,\infty)$, both integrands are non-negative, so the
    improper integrals are unambiguous (possibly $+\infty$). The volume is finite:
    \[\pi\int_1^\infty \frac{1}{x^2}\,dx = \pi\Big[-\frac1x\Big]_1^\infty = \boxed{\ \pi\ }.\]
    For the surface area, $f'(x) = -1/x^2$, so the integrand is $\sqrt{1+x^{-4}}/x \geq 1/x$
    pointwise, and since $\int_1^\infty \frac1x\,dx$ diverges, so does
    \[2\pi\int_1^\infty \frac{\sqrt{1+x^{-4}}}{x}\,dx = \boxed{\ +\infty\ }.\]
    This is the classical \newterm{Gabriel's horn} \germanfor{Torricellis Trompete}: finite
    volume, infinite surface area.

  \item ($\Leftarrow$) Suppose $(*)$ holds. Since $f'(a^+)=+\infty$, the inverse
    function theorem applies to $f|_{(a,a+\varepsilon)}$ for $\varepsilon$ small, giving a $C^1$,
    increasing inverse there; extend it by
    \[g(y) := \begin{cases} f^{-1}(y) & f(a)<y<f(a)+\varepsilon, \\ a & y\leq f(a). \end{cases}\]
    Since $f'(a^+)=+\infty$, $\lim_{y\to f(a)^+}g'(y) = \lim_{x\to a^+} 1/f'(x) = 0 =
    \lim_{y\to f(a)^-}g'(y)$, so $g\in C^1$ across $y=f(a)$. Around a boundary point $P :=
    (a,f(a)\cos\phi,f(a)\sin\phi)$, the map $(t,\theta)\mapsto\big(g(t),f(g(t))\cos\theta,
    f(g(t))\sin\theta\big)$ is then a regular $C^1$ parametrization of $\partial\Omega$ (using
    $f(a)>0$ to keep the angular direction non-degenerate). The symmetric argument at $b$, using
    $f'(b^-)=-\infty$, handles the other end, so $\partial\Omega$ is a $C^1$ submanifold and
    $\Omega$ is a bounded $C^1$ domain.

    ($\Rightarrow$) Suppose $\Omega$ is a bounded $C^1$ domain. If $f$ vanished at some interior
    point, $\Omega$ would be disconnected by part \textbf{(a)}, which
    \cref{def:bounded_ck_domain} excludes, since it requires $\Omega$ itself to make sense as a
    single bounded open set with $(n-1)$-manifold boundary (implicit in the exercise's setup).
    Therefore $f>0$ on $(a,b)$. For the endpoint condition at $a$: the point
    $P:=(a,f(a),0)\in\partial\Omega$ is a smooth boundary point, so
    $T_P(\partial\Omega)$ is $2$-dimensional and contains
    \[\frac{d}{dt}\Big|_{t=0^+}\big(a,f(a)-t,0\big) = -e_2, \qquad
      \frac{d}{dt}\Big|_{t=0^-}\big(a,f(a)\cos t,f(a)\sin t\big) = f(a)e_3,\]
    so $T_P(\partial\Omega) = \operatorname{span}\{e_2,e_3\}$. But $T_P(\partial\Omega)$ must
    also contain the limit of the (rescaled, if necessary) tangent vectors
    \[v_\varepsilon := \frac{d}{dt}\Big|_{t=0^+}\big(a+\varepsilon+t,\ f(a+\varepsilon+t),\
      0\big) = \big(1,\ f'(a+\varepsilon),\ 0\big)\]
    as $\varepsilon\to0^+$ (these are tangent to $\partial\Omega$ at nearby points converging to
    $P$). If $f'(a+\varepsilon)$ stayed bounded along some sequence $\varepsilon_k\to0^+$, the
    direction of $v_{\varepsilon_k}$ would converge to some $\alpha e_1+\beta e_2 \notin
    \operatorname{span}\{e_2,e_3\}$ (its $e_1$-component is the nonzero constant $1$, up to
    normalization), contradicting $T_P(\partial\Omega)=\operatorname{span}\{e_2,e_3\}$. So
    $f'(a^+)=+\infty$; the symmetric argument at $b$ gives $f'(b^-)=-\infty$.
\end{enumerate}
\end{exercisesolution}

\begin{ainote}
All parts of \cref{ex:spherical_quadrilateral_area,ex:torus_surface_area,ex:solids_of_revolution_area} were worked out independently before consulting
\texttt{Sol10\_Analysis2\_eng.pdf}. The three main answers ($(\varphi_2-\varphi_1)(\sin\theta_2
-\sin\theta_1)$, $8\pi^2$, $\pi\int_a^bf(x)^2dx$ / $2\pi\int_a^bf(x)\sqrt{1+f'(x)^2}dx$ /
volume $\pi$, area $+\infty$) match exactly. Part \textbf{(e)} was reconstructed independently
along the same two-directional strategy (build an explicit $C^1$ parametrization near each
endpoint using the inverse function theorem; for the converse, derive a contradiction from the
tangent space being forced to both contain $\operatorname{span}\{e_2,e_3\}$ and the limiting
direction of $v_\varepsilon$) and agrees with the official solution in substance, though the
write-up above is condensed relative to the official one.
\end{ainote}
